% =====================================================================
%  05  反応器設計 — 結論＋採用構成（What & Why を一枚で）
%  source_memo: 04_reaction.tex §4.1 反応器形式の選定の3判断と結論、
%               §4.x 検証諸元 Table reactor_verification（trial #201）
%  方針：上=3つの設計判断を課題→対策で3列／下左=採用構成の模式図(TikZ)／
%        下右=キー諸元。色は状態(反応/再生)が意味を持つ模式図のみ。丸括弧禁止。
% =====================================================================
\begin{frame}

  % ===== 反応の性質 → 採用部品 の対応マップ ========================
  \vspace{0.2em}
  \centering
  \begin{tikzpicture}[
    >=Stealth, font=\footnotesize,
    prop/.style={draw=black!50, rounded corners=2pt, align=center,
                 inner sep=3pt, minimum height=0.62cm, text width=2.6cm},
    part/.style={draw=HeaderBlue, fill=HiBlue!45, rounded corners=2pt, align=center,
                 inner sep=3pt, minimum height=0.5cm, text width=3.5cm,
                 font=\footnotesize\bfseries},
    link/.style={-{Stealth[length=2mm]}, ArrowBlue, line width=1pt},
  ]
    % ---- 左：反応の性質 ----
    \node[prop] (heat) at (0,2.35) {吸熱};
    \node[prop] (eq)   at (0,1.15) {平衡支配\\高温・低圧 $0.5$ bar};
    \node[prop] (coke) at (0,0)    {コーキング\\数分で失活};
    % ---- 右：採用部品 ----
    \node[part] (hgm)   at (5.4,2.4) {床内 HGM 補温};
    \node[part] (shal)  at (5.4,1.6) {浅床軸流};
    \node[part] (multi) at (5.4,0.8) {多基・大断面・低速};
    \node[part] (swing) at (5.4,0.0) {スイング　反応 $\rightleftarrows$ 再生};
    % ---- 結線（どの性質がどの部品を強制するか）----
    \draw[link] (heat.east) -- (hgm.west);
    \draw[link] (eq.east)   -- (shal.west);
    \draw[link] (eq.east)   -- (multi.west);
    \draw[link] (coke.east) -- (swing.west);
    % ---- 見出し ----
    \node[font=\footnotesize\bfseries, anchor=south] at (0,2.85)   {反応の性質};
    \node[font=\footnotesize\bfseries, anchor=south] at (5.4,2.95) {採用部品};
  \end{tikzpicture}

  \vspace{0.05em}

  % ===== 下段：採用構成の模式図（左）／キー諸元（右）===============
  \begin{columns}[T,onlytextwidth]
    \begin{column}{0.58\textwidth}
      {\bfseries 採用構成　浅床軸流・多基並列スイング}\par
      \vspace{0.1em}
      \centering
      \scalebox{0.72}{%
      \begin{tikzpicture}[
        >=Stealth, font=\scriptsize,
        wall/.style={line width=1pt},
        proc/.style={-{Stealth[length=2mm]}, line width=1pt},
        flow/.style={-{Stealth[length=1.4mm]}, line width=0.5pt, black!65},
        dim/.style={{Stealth[length=1.6mm]}-{Stealth[length=1.6mm]}, line width=0.4pt},
        lead/.style={line width=0.3pt, black!55},
        lbl/.style={font=\scriptsize, align=left},
      ]
        \def\R{1.8}\def\Hb{1.35}\def\ry{0.62}
        % ---- 短胴・球形鏡板（扁平な大径容器）----
        \draw[wall] (-\R,0) -- (-\R,\Hb);
        \draw[wall] (\R,0) -- (\R,\Hb);
        \draw[wall] (-\R,\Hb) arc[start angle=180, end angle=0, x radius=\R, y radius=\ry];
        \draw[wall] (\R,0) arc[start angle=0, end angle=-180, x radius=\R, y radius=\ry];
        % ---- 入口ノズル（上）----
        \draw[wall] (-0.26,\Hb+\ry-0.02) -- (-0.26,\Hb+\ry+0.4);
        \draw[wall] (0.26,\Hb+\ry-0.02) -- (0.26,\Hb+\ry+0.4);
        \draw[wall] (-0.38,\Hb+\ry+0.4) rectangle (0.38,\Hb+\ry+0.5);
        \draw[proc] (0,\Hb+\ry+1.0) -- (0,\Hb+\ry+0.55);
        \node[lbl, anchor=west] at (0.45,\Hb+\ry+0.8) {原料ガス};
        % ---- 出口ノズル（下）----
        \draw[wall] (-0.26,-\ry+0.02) -- (-0.26,-\ry-0.4);
        \draw[wall] (0.26,-\ry+0.02) -- (0.26,-\ry-0.4);
        \draw[wall] (-0.38,-\ry-0.5) rectangle (0.38,-\ry-0.4);
        \draw[proc] (0,-\ry-0.55) -- (0,-\ry-1.0);
        \node[lbl, anchor=west] at (0.45,-\ry-0.8) {生成ガス};
        % ---- 分配板（多孔板）----
        \draw[line width=0.8pt] (-\R+0.05,1.08) -- (\R-0.05,1.08);
        \foreach \x in {-1.4,-1.0,-0.6,-0.2,0.2,0.6,1.0,1.4}{\draw[line width=0.4pt] (\x,1.08) -- (\x,0.98);}
        % ---- 触媒・HGM 混合浅床 ----
        \fill[HiBlue] (-\R+0.05,0.6) rectangle (\R-0.05,0.93);
        \draw[black!45] (-\R+0.05,0.6) rectangle (\R-0.05,0.93);
        % ---- 支持グリッド ----
        \draw[line width=0.8pt] (-\R+0.05,0.55) -- (\R-0.05,0.55);
        % ---- 床内 軸流（下向き）----
        \foreach \x in {-1.1,0,1.1}{\draw[flow] (\x,0.96) -- (\x,0.57);}
        % ---- 左ラベル ----
        \node[lbl, anchor=east] (ld) at (-\R-0.45,1.08) {分配板}; \draw[lead] (ld.east) -- (-\R+0.05,1.08);
        \node[lbl, anchor=east] (lc) at (-\R-0.45,0.77) {触媒・HGM 床}; \draw[lead] (lc.east) -- (-\R+0.05,0.77);
        \node[lbl, anchor=east] (lg) at (-\R-0.45,0.5) {支持グリッド}; \draw[lead] (lg.east) -- (-\R+0.05,0.55);
        % ---- 寸法 ----
        \draw[dim] (-\R,-1.2) -- (\R,-1.2); \node[below] at (0,-1.22) {$D$};
        \draw[dim] (\R+0.32,0.6) -- (\R+0.32,0.93); \node[right] at (\R+0.38,0.765) {$L_{\mathrm{bed}}$};
      \end{tikzpicture}}\par
      \vspace{0.1em}
      {\scriptsize 1 基の断面　$L_{\mathrm{bed}}\!\ll\! D$、常時 $7$／全 $14$ 基}\par
    \end{column}
    \begin{column}{0.42\textwidth}
      % 諸元は最良設計点 trial #201（スライドには番号を出さない）
      {\bfseries キー諸元}\par
      \vspace{0.4em}
      {\setlength{\fboxsep}{5pt}\setlength{\fboxrule}{1pt}%
       \fcolorbox{HeaderBlue}{white}{%
        \footnotesize\renewcommand{\arraystretch}{1.08}%
        \begin{tabular}{@{}l@{\hspace{1em}}l@{}}
          総基数 $N_{\mathrm{total}}$        & $14$ 基 \\
          並列基数 $N_{\mathrm{online}}$      & $7$ \\
          浅床厚 $L_{\mathrm{bed}}$           & $1.58$ m \\
          触媒粒径 $d_p$                      & $5.84$ mm \\
          入口温度 $T_{\mathrm{in}}$          & $935$ K \\
          床温降下 $\Delta T_{\max}$          & $50$ K \\
          有効触媒分率 $\phi_{\mathrm{cat}}$  & $0.85$ \\
          HGM 補償熱 $Q_{\mathrm{HGM}}$       & $24.96$ MW \\
        \end{tabular}}}
    \end{column}
  \end{columns}

\end{frame}
